\documentclass[12pt,a4paper]{article}
\usepackage[utf8]{vietnam}
\usepackage{amsmath,amssymb}
\usepackage{tabularx}
\usepackage{tikz}
\usepackage{xcolor}
\usepackage[top=0.6cm, bottom=0.6cm, left=0.6cm, right=0.6cm]{geometry}
\usepackage[dethi]{ex_test}

\SAOPTN{widthSA=3.2,heightSA=0.62}

\setlength{\parskip}{0pt}

% Dòng kẻ chấm mờ đều đặn chuẩn viết tay
\newcommand{\dongke}[1]{%
  \begingroup%
  \setlength{\parskip}{0pt}%
  \baselineskip=18.5pt%
  \vspace{1.0mm}%
  \foreach \i in {1,...,#1}{%
    \noindent{\color{gray!60}\dotfill}\par%
  }%
  \vspace{0.6mm}%
  \endgroup%
}

% Định nghĩa ô bài tập tự luận trong bảng
\newcommand{\oBaiTap}[3][6]{%
  \vspace{1.0mm}%
  \noindent\textbf{#2} #3\par%
  \dongke{#1}%
}

\begin{document}
\pagestyle{empty}
% =============================================================================
% ========================= TOÁN 9 - TUẦN 07 - BUỔI 01: MÃ ĐỀ 101 =============
% ========================= MẶT TRƯỚC (TRANG 1/2) ==============================
% =============================================================================

\noindent
\begin{tabularx}{\linewidth}{|X|c|}
\hline
\textbf{PHIẾU KIỂM TRA ĐỊNH KỲ TOÁN 9 - TUẦN 07 - BUỔI 01} & \textbf{Mã đề: 101 - Trang 1/2} \\
\hline
\end{tabularx}

\vspace{1.0mm}

\noindent
\begin{tabularx}{\linewidth}{|p{3.2cm}|X|}
\hline
\centering\textbf{ĐIỂM SỐ} & \centering\textbf{LỜI NHẬN XÉT CỦA THẦY CÔ} \tabularnewline
\hline
\rule{0pt}{1.8cm} & \tabularnewline
\hline
\end{tabularx}

\vspace{1.2mm}
\noindent\textbf{Câu 1 (0,25 điểm).} Cho $\Delta ABC$ vuông tại $A$ có $AB = 6, BC = 10$. Khi đó $\sin C =$ \shortans{\dfrac{3}{5}}\par\vspace{0.8mm}
\noindent\textbf{Câu 2 (0,25 điểm).} Cho $\Delta ABC$ vuông tại $A$ có $AC = 8, BC = 10$. Khi đó $\cos C =$ \shortans{\dfrac{4}{5}}\par\vspace{0.8mm}
\noindent\textbf{Câu 3 (0,25 điểm).} Cho $\Delta ABC$ vuông tại $A$ có $AB = 3, AC = 4$. Khi đó $\tan B =$ \shortans{\dfrac{4}{3}}\par\vspace{0.8mm}
\noindent\textbf{Câu 4 (0,25 điểm).} Cho góc nhọn $\alpha$ có $\sin \alpha = 0{,}75$. Khi đó $\cos(90^\circ - \alpha) =$ \shortans{0{,}75}\par\vspace{0.8mm}
\noindent\textbf{Câu 5 (0,25 điểm).} Giá trị của biểu thức $\sin^2 25^\circ + \sin^2 65^\circ$ bằng: \shortans{1}\par\vspace{0.8mm}
\noindent\textbf{Câu 6 (0,25 điểm).} Giá trị của biểu thức $\tan 45^\circ - 2\sin 30^\circ$ bằng: \shortans{0}\par\vspace{0.8mm}
\noindent\textbf{Câu 7 (0,25 điểm).} Cột cờ cao $7\text{ m}$ có bóng trên mặt đất dài $7\text{ m}$. Góc của tia nắng là: \shortans{45^\circ}\par\vspace{0.8mm}
\noindent\textbf{Câu 8 (0,25 điểm).} Cặp $(1; y)$ là nghiệm của phương trình $3x + y = 7$. Khi đó $y =$ \shortans{4}\par\vspace{0.8mm}
\noindent\textbf{Câu 9 (0,25 điểm).} Tìm $m$ để cặp $(2; 1)$ là nghiệm của phương trình $mx + 3y = 7$: \shortans{2}\par\vspace{0.8mm}
\noindent\textbf{Câu 10 (0,25 điểm).} Nghiệm dương của phương trình $(x - 4)(x + 1) = 0$ là: \shortans{4}\par\vspace{0.8mm}
\vspace{0.8mm}
\noindent\textbf{Câu 11 (2,5 điểm). Đại số tổng hợp:}
\vspace{0.6mm}

\noindent
\begin{tabularx}{\linewidth}{|X|X|}
\hline
\oBaiTap[22]{a) (1,25 điểm)}{$\begin{cases} 2x + y = 8 \\ 3x - y = 7 \end{cases}$} & 
\oBaiTap[22]{b) (1,25 điểm)}{Rút gọn với $x > 0, x \ne 4$:\par $P = \left(\dfrac{\sqrt{x}}{\sqrt{x}-2} - \dfrac{4}{x-2\sqrt{x}}\right) : \dfrac{\sqrt{x}+2}{\sqrt{x}}$} \\
\hline
\end{tabularx}

\newpage
% =============================================================================
% ========================= TOÁN 9 - TUẦN 07 - BUỔI 01: MÃ ĐỀ 101 =============
% ========================= MẶT SAU (TRANG 2/2) ================================
% =============================================================================

\noindent
\begin{tabularx}{\linewidth}{|X|c|}
\hline
\textbf{TOÁN 9 - TUẦN 07 - BUỔI 01} & \textbf{Mã đề: 101 - Trang 2/2} \\
\hline
\end{tabularx}

\vspace{1.2mm}

\noindent
\begin{tabularx}{\linewidth}{|X|}
\hline
\vspace{0.5mm}
\textbf{Câu 12 (2,5 điểm).} Hai bạn An và Bình cùng đến nhà sách mua đồ dùng học tập chuẩn bị cho năm học mới. Bạn An mua $3$ quyển vở và $2$ chiếc bút bi cùng loại hết tổng cộng $34\,000$ đồng. Bạn Bình mua $2$ quyển vở và $4$ chiếc bút bi cùng loại hết tổng cộng $36\,000$ đồng. Tính giá tiền của mỗi quyển vở và mỗi chiếc bút bi.\par
\vspace{0.4mm}
\noindent\textit{Bài giải:}\dongke{16}
\\ \hline
\end{tabularx}

\vspace{2.5mm}

\noindent
\begin{tabularx}{\linewidth}{|X|}
\hline
\vspace{0.5mm}
\textbf{Câu 13 (2,5 điểm).} Cho tam giác $ABC$ vuông tại $A$, đường cao $AH$. Biết $AB = 6\text{ cm}, BC = 10\text{ cm}$.\par
\textbf{a) (1,5 điểm)} Tính độ dài cạnh $AC$, số đo góc $B$ (làm tròn đến độ) và độ dài đường cao $AH$.\par
\textbf{b) (1,0 điểm)} Kẻ $HE \perp AB$ tại $E$ và $HF \perp AC$ tại $F$. Chứng minh rằng: $AE \cdot AB = AF \cdot AC$.\par
\vspace{0.4mm}
\noindent\textit{Bài giải:}\dongke{17}
\\ \hline
\end{tabularx}

\newpage
% =============================================================================
% ========================= TOÁN 9 - TUẦN 07 - BUỔI 01: MÃ ĐỀ 102 =============
% ========================= MẶT TRƯỚC (TRANG 1/2) ==============================
% =============================================================================

\noindent
\begin{tabularx}{\linewidth}{|X|c|}
\hline
\textbf{PHIẾU KIỂM TRA ĐỊNH KỲ TOÁN 9 - TUẦN 07 - BUỔI 01} & \textbf{Mã đề: 102 - Trang 1/2} \\
\hline
\end{tabularx}

\vspace{1.0mm}

\noindent
\begin{tabularx}{\linewidth}{|p{3.2cm}|X|}
\hline
\centering\textbf{ĐIỂM SỐ} & \centering\textbf{LỜI NHẬN XÉT CỦA THẦY CÔ} \tabularnewline
\hline
\rule{0pt}{1.8cm} & \tabularnewline
\hline
\end{tabularx}

\vspace{1.2mm}
\noindent\textbf{Câu 1 (0,25 điểm).} Cho $\Delta ABC$ vuông tại $A$ có $AB = 8, BC = 10$. Khi đó $\sin C =$ \shortans{\dfrac{4}{5}}\par\vspace{0.8mm}
\noindent\textbf{Câu 2 (0,25 điểm).} Cho $\Delta ABC$ vuông tại $A$ có $AC = 6, BC = 10$. Khi đó $\cos C =$ \shortans{\dfrac{3}{5}}\par\vspace{0.8mm}
\noindent\textbf{Câu 3 (0,25 điểm).} Cho $\Delta ABC$ vuông tại $A$ có $AB = 4, AC = 3$. Khi đó $\tan B =$ \shortans{\dfrac{3}{4}}\par\vspace{0.8mm}
\noindent\textbf{Câu 4 (0,25 điểm).} Cho góc nhọn $\alpha$ có $\cos \alpha = 0{,}65$. Khi đó $\sin(90^\circ - \alpha) =$ \shortans{0{,}65}\par\vspace{0.8mm}
\noindent\textbf{Câu 5 (0,25 điểm).} Giá trị của biểu thức $\cos^2 35^\circ + \cos^2 55^\circ$ bằng: \shortans{1}\par\vspace{0.8mm}
\noindent\textbf{Câu 6 (0,25 điểm).} Giá trị của biểu thức $2\cos 60^\circ - \cot 45^\circ$ bằng: \shortans{0}\par\vspace{0.8mm}
\noindent\textbf{Câu 7 (0,25 điểm).} Bóng cây dài $12\text{ m}$, tia nắng tạo góc $45^\circ$ với đất. Chiều cao cây là: \shortans{12\text{ m}}\par\vspace{0.8mm}
\noindent\textbf{Câu 8 (0,25 điểm).} Cặp $(2; y)$ là nghiệm của phương trình $2x + y = 9$. Khi đó $y =$ \shortans{5}\par\vspace{0.8mm}
\noindent\textbf{Câu 9 (0,25 điểm).} Tìm $m$ để cặp $(1; 3)$ là nghiệm của phương trình $mx + 2y = 8$: \shortans{2}\par\vspace{0.8mm}
\noindent\textbf{Câu 10 (0,25 điểm).} Nghiệm dương của phương trình $(x - 5)(x + 3) = 0$ là: \shortans{5}\par\vspace{0.8mm}
\vspace{0.8mm}
\noindent\textbf{Câu 11 (2,5 điểm). Đại số tổng hợp:}
\vspace{0.6mm}

\noindent
\begin{tabularx}{\linewidth}{|X|X|}
\hline
\oBaiTap[22]{a) (1,25 điểm)}{$\begin{cases} 3x - y = 5 \\ x + y = 7 \end{cases}$} & 
\oBaiTap[22]{b) (1,25 điểm)}{Rút gọn với $x > 0, x \ne 9$:\par $P = \left(\dfrac{\sqrt{x}}{\sqrt{x}-3} - \dfrac{9}{x-3\sqrt{x}}\right) : \dfrac{\sqrt{x}+3}{\sqrt{x}}$} \\
\hline
\end{tabularx}

\newpage
% =============================================================================
% ========================= TOÁN 9 - TUẦN 07 - BUỔI 01: MÃ ĐỀ 102 =============
% ========================= MẶT SAU (TRANG 2/2) ================================
% =============================================================================

\noindent
\begin{tabularx}{\linewidth}{|X|c|}
\hline
\textbf{TOÁN 9 - TUẦN 07 - BUỔI 01} & \textbf{Mã đề: 102 - Trang 2/2} \\
\hline
\end{tabularx}

\vspace{1.2mm}

\noindent
\begin{tabularx}{\linewidth}{|X|}
\hline
\vspace{0.5mm}
\textbf{Câu 12 (2,5 điểm).} Hai bạn An và Bình cùng đến nhà sách mua đồ dùng học tập chuẩn bị cho năm học mới. Bạn An mua $4$ quyển vở và $2$ chiếc bút bi cùng loại hết tổng cộng $42\,000$ đồng. Bạn Bình mua $3$ quyển vở và $5$ chiếc bút bi cùng loại hết tổng cộng $49\,000$ đồng. Tính giá tiền của mỗi quyển vở và mỗi chiếc bút bi.\par
\vspace{0.4mm}
\noindent\textit{Bài giải:}\dongke{16}
\\ \hline
\end{tabularx}

\vspace{2.5mm}

\noindent
\begin{tabularx}{\linewidth}{|X|}
\hline
\vspace{0.5mm}
\textbf{Câu 13 (2,5 điểm).} Cho tam giác $ABC$ vuông tại $A$, đường cao $AH$. Biết $AC = 8\text{ cm}, BC = 10\text{ cm}$.\par
\textbf{a) (1,5 điểm)} Tính độ dài cạnh $AB$, số đo góc $C$ (làm tròn đến độ) và độ dài đường cao $AH$.\par
\textbf{b) (1,0 điểm)} Kẻ $HE \perp AB$ tại $E$ và $HF \perp AC$ tại $F$. Chứng minh tứ giác $AEHF$ là hình chữ nhật và suy ra $EF = AH$.\par
\vspace{0.4mm}
\noindent\textit{Bài giải:}\dongke{17}
\\ \hline
\end{tabularx}

\newpage
% =============================================================================
% ========================= TOÁN 9 - TUẦN 07 - BUỔI 01: MÃ ĐỀ 103 =============
% ========================= MẶT TRƯỚC (TRANG 1/2) ==============================
% =============================================================================

\noindent
\begin{tabularx}{\linewidth}{|X|c|}
\hline
\textbf{PHIẾU KIỂM TRA ĐỊNH KỲ TOÁN 9 - TUẦN 07 - BUỔI 01} & \textbf{Mã đề: 103 - Trang 1/2} \\
\hline
\end{tabularx}

\vspace{1.0mm}

\noindent
\begin{tabularx}{\linewidth}{|p{3.2cm}|X|}
\hline
\centering\textbf{ĐIỂM SỐ} & \centering\textbf{LỜI NHẬN XÉT CỦA THẦY CÔ} \tabularnewline
\hline
\rule{0pt}{1.8cm} & \tabularnewline
\hline
\end{tabularx}

\vspace{1.2mm}
\noindent\textbf{Câu 1 (0,25 điểm).} Cho $\Delta ABC$ vuông tại $A$ có $AB = 5, BC = 13$. Khi đó $\sin C =$ \shortans{\dfrac{5}{13}}\par\vspace{0.8mm}
\noindent\textbf{Câu 2 (0,25 điểm).} Cho $\Delta ABC$ vuông tại $A$ có $AC = 12, BC = 13$. Khi đó $\cos C =$ \shortans{\dfrac{12}{13}}\par\vspace{0.8mm}
\noindent\textbf{Câu 3 (0,25 điểm).} Cho $\Delta ABC$ vuông tại $A$ có $AB = 5, AC = 12$. Khi đó $\tan B =$ \shortans{\dfrac{12}{5}}\par\vspace{0.8mm}
\noindent\textbf{Câu 4 (0,25 điểm).} Cho góc nhọn $\alpha$ có $\sin \alpha = \dfrac{5}{13}$. Khi đó $\cos(90^\circ - \alpha) =$ \shortans{\dfrac{5}{13}}\par\vspace{0.8mm}
\noindent\textbf{Câu 5 (0,25 điểm).} Giá trị của biểu thức $\sin^2 40^\circ + \sin^2 50^\circ$ bằng: \shortans{1}\par\vspace{0.8mm}
\noindent\textbf{Câu 6 (0,25 điểm).} Giá trị của biểu thức $\sqrt{3}\tan 30^\circ - 1$ bằng: \shortans{0}\par\vspace{0.8mm}
\noindent\textbf{Câu 7 (0,25 điểm).} Cách chân hải đăng $25\text{ m}$, góc nhìn lên đỉnh là $45^\circ$. Hải đăng cao: \shortans{25\text{ m}}\par\vspace{0.8mm}
\noindent\textbf{Câu 8 (0,25 điểm).} Cặp $(3; y)$ là nghiệm của phương trình $x + 2y = 7$. Khi đó $y =$ \shortans{2}\par\vspace{0.8mm}
\noindent\textbf{Câu 9 (0,25 điểm).} Tìm $m$ để cặp $(3; 1)$ là nghiệm của phương trình $mx + 4y = 10$: \shortans{2}\par\vspace{0.8mm}
\noindent\textbf{Câu 10 (0,25 điểm).} Nghiệm dương của phương trình $(x - 7)(x + 2) = 0$ là: \shortans{7}\par\vspace{0.8mm}
\vspace{0.8mm}
\noindent\textbf{Câu 11 (2,5 điểm). Đại số tổng hợp:}
\vspace{0.6mm}

\noindent
\begin{tabularx}{\linewidth}{|X|X|}
\hline
\oBaiTap[22]{a) (1,25 điểm)}{$\begin{cases} 2x + y = 9 \\ x - y = 3 \end{cases}$} & 
\oBaiTap[22]{b) (1,25 điểm)}{Rút gọn với $x > 0, x \ne 1$:\par $P = \left(\dfrac{\sqrt{x}}{\sqrt{x}-1} - \dfrac{1}{x-\sqrt{x}}\right) : \dfrac{\sqrt{x}+1}{\sqrt{x}}$} \\
\hline
\end{tabularx}

\newpage
% =============================================================================
% ========================= TOÁN 9 - TUẦN 07 - BUỔI 01: MÃ ĐỀ 103 =============
% ========================= MẶT SAU (TRANG 2/2) ================================
% =============================================================================

\noindent
\begin{tabularx}{\linewidth}{|X|c|}
\hline
\textbf{TOÁN 9 - TUẦN 07 - BUỔI 01} & \textbf{Mã đề: 103 - Trang 2/2} \\
\hline
\end{tabularx}

\vspace{1.2mm}

\noindent
\begin{tabularx}{\linewidth}{|X|}
\hline
\vspace{0.5mm}
\textbf{Câu 12 (2,5 điểm).} Hai bạn An và Bình cùng đến nhà sách mua đồ dùng học tập chuẩn bị cho năm học mới. Bạn An mua $2$ quyển vở và $3$ chiếc bút bi cùng loại hết tổng cộng $31\,000$ đồng. Bạn Bình mua $5$ quyển vở và $2$ chiếc bút bi cùng loại hết tổng cộng $50\,000$ đồng. Tính giá tiền của mỗi quyển vở và mỗi chiếc bút bi.\par
\vspace{0.4mm}
\noindent\textit{Bài giải:}\dongke{16}
\\ \hline
\end{tabularx}

\vspace{2.5mm}

\noindent
\begin{tabularx}{\linewidth}{|X|}
\hline
\vspace{0.5mm}
\textbf{Câu 13 (2,5 điểm).} Cho tam giác $ABC$ vuông tại $A$, đường cao $AH$. Biết $AB = 5\text{ cm}, BC = 13\text{ cm}$.\par
\textbf{a) (1,5 điểm)} Tính độ dài cạnh $AC$, số đo góc $B$ (làm tròn đến độ) và độ dài đường cao $AH$ (làm tròn đến chữ số thập phân thứ nhất).\par
\textbf{b) (1,0 điểm)} Kẻ $HE \perp AB$ tại $E$ và $HF \perp AC$ tại $F$. Chứng minh rằng: $AE \cdot AB = AF \cdot AC$.\par
\vspace{0.4mm}
\noindent\textit{Bài giải:}\dongke{17}
\\ \hline
\end{tabularx}

\newpage
% =============================================================================
% ========================= TOÁN 9 - TUẦN 07 - BUỔI 01: MÃ ĐỀ 104 =============
% ========================= MẶT TRƯỚC (TRANG 1/2) ==============================
% =============================================================================

\noindent
\begin{tabularx}{\linewidth}{|X|c|}
\hline
\textbf{PHIẾU KIỂM TRA ĐỊNH KỲ TOÁN 9 - TUẦN 07 - BUỔI 01} & \textbf{Mã đề: 104 - Trang 1/2} \\
\hline
\end{tabularx}

\vspace{1.0mm}

\noindent
\begin{tabularx}{\linewidth}{|p{3.2cm}|X|}
\hline
\centering\textbf{ĐIỂM SỐ} & \centering\textbf{LỜI NHẬN XÉT CỦA THẦY CÔ} \tabularnewline
\hline
\rule{0pt}{1.8cm} & \tabularnewline
\hline
\end{tabularx}

\vspace{1.2mm}
\noindent\textbf{Câu 1 (0,25 điểm).} Cho $\Delta ABC$ vuông tại $A$ có $AB = 12, BC = 13$. Khi đó $\sin C =$ \shortans{\dfrac{12}{13}}\par\vspace{0.8mm}
\noindent\textbf{Câu 2 (0,25 điểm).} Cho $\Delta ABC$ vuông tại $A$ có $AC = 5, BC = 13$. Khi đó $\cos C =$ \shortans{\dfrac{5}{13}}\par\vspace{0.8mm}
\noindent\textbf{Câu 3 (0,25 điểm).} Cho $\Delta ABC$ vuông tại $A$ có $AB = 12, AC = 5$. Khi đó $\tan B =$ \shortans{\dfrac{5}{12}}\par\vspace{0.8mm}
\noindent\textbf{Câu 4 (0,25 điểm).} Cho góc nhọn $\alpha$ có $\cos \alpha = \dfrac{7}{25}$. Khi đó $\sin(90^\circ - \alpha) =$ \shortans{\dfrac{7}{25}}\par\vspace{0.8mm}
\noindent\textbf{Câu 5 (0,25 điểm).} Giá trị của biểu thức $\cos^2 20^\circ + \cos^2 70^\circ$ bằng: \shortans{1}\par\vspace{0.8mm}
\noindent\textbf{Câu 6 (0,25 điểm).} Giá trị của biểu thức $\tan 45^\circ - \sqrt{3}\cot 60^\circ$ bằng: \shortans{0}\par\vspace{0.8mm}
\noindent\textbf{Câu 7 (0,25 điểm).} Thang dài $6\text{ m}$ dựa tường tạo góc $60^\circ$. Chân thang cách tường: \shortans{3\text{ m}}\par\vspace{0.8mm}
\noindent\textbf{Câu 8 (0,25 điểm).} Cặp $(1; y)$ là nghiệm của phương trình $4x + y = 10$. Khi đó $y =$ \shortans{6}\par\vspace{0.8mm}
\noindent\textbf{Câu 9 (0,25 điểm).} Tìm $m$ để cặp $(2; 2)$ là nghiệm của phương trình $mx + y = 8$: \shortans{3}\par\vspace{0.8mm}
\noindent\textbf{Câu 10 (0,25 điểm).} Nghiệm dương của phương trình $(x - 6)(x + 4) = 0$ là: \shortans{6}\par\vspace{0.8mm}
\vspace{0.8mm}
\noindent\textbf{Câu 11 (2,5 điểm). Đại số tổng hợp:}
\vspace{0.6mm}

\noindent
\begin{tabularx}{\linewidth}{|X|X|}
\hline
\oBaiTap[22]{a) (1,25 điểm)}{$\begin{cases} 3x + y = 11 \\ 2x - y = 4 \end{cases}$} & 
\oBaiTap[22]{b) (1,25 điểm)}{Rút gọn với $x > 0, x \ne 16$:\par $P = \left(\dfrac{\sqrt{x}}{\sqrt{x}-4} - \dfrac{16}{x-4\sqrt{x}}\right) : \dfrac{\sqrt{x}+4}{\sqrt{x}}$} \\
\hline
\end{tabularx}

\newpage
% =============================================================================
% ========================= TOÁN 9 - TUẦN 07 - BUỔI 01: MÃ ĐỀ 104 =============
% ========================= MẶT SAU (TRANG 2/2) ================================
% =============================================================================

\noindent
\begin{tabularx}{\linewidth}{|X|c|}
\hline
\textbf{TOÁN 9 - TUẦN 07 - BUỔI 01} & \textbf{Mã đề: 104 - Trang 2/2} \\
\hline
\end{tabularx}

\vspace{1.2mm}

\noindent
\begin{tabularx}{\linewidth}{|X|}
\hline
\vspace{0.5mm}
\textbf{Câu 12 (2,5 điểm).} Hai bạn An và Bình cùng đến nhà sách mua đồ dùng học tập chuẩn bị cho năm học mới. Bạn An mua $3$ quyển vở và $4$ chiếc bút bi cùng loại hết tổng cộng $44\,000$ đồng. Bạn Bình mua $4$ quyển vở và $2$ chiếc bút bi cùng loại hết tổng cộng $42\,000$ đồng. Tính giá tiền của mỗi quyển vở và mỗi chiếc bút bi.\par
\vspace{0.4mm}
\noindent\textit{Bài giải:}\dongke{16}
\\ \hline
\end{tabularx}

\vspace{2.5mm}

\noindent
\begin{tabularx}{\linewidth}{|X|}
\hline
\vspace{0.5mm}
\textbf{Câu 13 (2,5 điểm).} Cho tam giác $ABC$ vuông tại $A$, đường cao $AH$. Biết $AC = 12\text{ cm}, BC = 13\text{ cm}$.\par
\textbf{a) (1,5 điểm)} Tính độ dài cạnh $AB$, số đo góc $C$ (làm tròn đến độ) và độ dài đường cao $AH$ (làm tròn đến chữ số thập phân thứ nhất).\par
\textbf{b) (1,0 điểm)} Kẻ $HE \perp AB$ tại $E$ và $HF \perp AC$ tại $F$. Chứng minh tứ giác $AEHF$ là hình chữ nhật và suy ra $EF = AH$.\par
\vspace{0.4mm}
\noindent\textit{Bài giải:}\dongke{17}
\\ \hline
\end{tabularx}

\end{document}
